\begin{table*}[!tbp]\centering
\caption{Frozen solver configurations. The three labelled variants share the same hierarchy implementation.}\label{tab:configuration}
\begin{tabular}{llll}
\toprule
Method & Fine labels / aggregation & Fine relaxation & Coarsest limit\\
\midrule
ML-CPMG & Persistent watershed, $B_{\max}=256$ & Vertical line on SPE10 &512\\
Min-threshold & Connected foreground $k>\min k$ & Same as ML-CPMG &512\\
Single-median & Connected foreground $\log k>\med(\log k)$ & Same as ML-CPMG &512\\
Algebraic-only & Heavy-edge aggregation from level 0 & Same as ML-CPMG &512\\
SA & Symmetric strength, $\theta=0.04$ & Symmetric Gauss--Seidel &10\\
RS & Classical strength, $\theta=0.25$ & Symmetric Gauss--Seidel &512\\
\bottomrule
\end{tabular}
\par\smallskip\footnotesize
Internal methods use $\omega_P=0.45$, $\tau_P=0$ on the fine level and $0.04$
thereafter, full positive-edge lumping, $\eta_{\rm stall}=0.33$, one pre- and
one post-smoothing step, and at most eight nodes per algebraic aggregate.
Synthetic grids use the common automatic line trigger; Brugge uses point
relaxation. Persistence, barrier, background and pedestal parameters are
$16$, $8$, $0.22$ and $2$, respectively. SA thresholds $0$ and $0.1$ are
reported in the sensitivity study.\normalsize
\end{table*}
